\chapter{Project Work}

Questo capitolo descrive il lavoro progettuale svolto durante la tesi.
In particolare, vengono presentati il contesto applicativo, i requisiti
del sistema, le scelte architetturali, le attività di progettazione e
implementazione della soluzione proposta.

Nota: questo capitolo deve descrivere il lavoro progettuale svolto
dallo studente e non deve contenere contenuti teorici già discussi
nel capitolo di stato dell’arte.

\section{Scopo del Progetto}

In questa sezione deve essere descritto il problema affrontato e gli
obiettivi del progetto.

Lo studente dovrebbe chiarire:

\begin{itemize}
\item il problema che il sistema intende risolvere;
\item le motivazioni del progetto;
\item gli obiettivi principali della soluzione proposta;
\item eventuali vincoli tecnici o applicativi.
\end{itemize}

\subsection{Contesto e Ambito di Applicazione}

Descrivere il contesto in cui si colloca il progetto, ad esempio:

\begin{itemize}
\item il dominio applicativo;
\item eventuali sistemi esistenti o infrastrutture di riferimento;
\item gli utenti o stakeholder coinvolti.
\end{itemize}

\subsection{Specifica dei Requisiti}

Descrivere i requisiti del sistema sviluppato.

\subsubsection{Requisiti Funzionali}

Elencare e descrivere le principali funzionalità che il sistema deve
fornire agli utenti.

\subsubsection{Requisiti Non Funzionali}

Descrivere eventuali requisiti relativi a:

\begin{itemize}
\item prestazioni;
\item sicurezza;
\item scalabilità;
\item affidabilità;
\item usabilità.
\end{itemize}

\subsection{Architettura di Sistema}

Descrivere l’architettura generale della soluzione sviluppata.

La sezione dovrebbe includere:

\begin{itemize}
\item descrizione dei principali componenti del sistema;
\item interazioni tra i moduli;
\item eventuali diagrammi architetturali.
\end{itemize}

\subsection{Progettazione}

Descrivere le scelte progettuali adottate per la realizzazione del sistema.

È possibile includere:

\begin{itemize}
\item diagrammi UML;
\item struttura dei moduli software;
\item progettazione delle API;
\item modello dei dati o schema del database.
\end{itemize}

\subsection{Implementazione e Tecniche di Sviluppo}

Descrivere le tecnologie utilizzate e le principali scelte implementative.

Ad esempio:

\begin{itemize}
\item linguaggi di programmazione;
\item framework utilizzati;
\item librerie principali;
\item strumenti di sviluppo.
\end{itemize}

\section{Testing e Validazione}

In questa sezione devono essere descritte le attività di testing
utilizzate per verificare il corretto funzionamento del sistema.

\subsection{Strategia di Testing}

Descrivere la strategia di testing adottata (ad esempio unit testing,
integration testing, system testing).

\subsection{Test Case}

Presentare alcuni esempi di test case significativi utilizzati
per verificare le funzionalità principali del sistema.

\subsection{Risultati dei Test}

Descrivere i risultati ottenuti dai test e discutere eventuali
problemi individuati e le relative soluzioni.

\section{Discussione dei Risultati}

Analizzare i risultati ottenuti e discutere:

\begin{itemize}
\item i punti di forza della soluzione proposta;
\item eventuali limitazioni del sistema;
\item possibili sviluppi futuri.
\end{itemize}